%======================================================================
\section{Overview and Roadmap}
\label{sec:overview}
%======================================================================

This memo collects four short technical notes. Each can be read
independently, but they were drafted to be read in sequence: the first
frames IHS structurally, the second tests how far the framing carries, the
third broadens the mapping to other OR paradigms, and the fourth connects
it to empirical practice.

\begin{itemize}
\item \textbf{\Cref{sec:part1} (Decomposition View):} develops the decomposition view. After reviewing the
core-guided and IHS algorithm families and the LBBD and Dantzig--Wolfe
templates, it identifies the row--column duality between cores and cuts,
gives a small worked example, and closes with a catalog of OR techniques
whose transfer to MaxSAT is supported by the analogy.

\item \textbf{\Cref{sec:part2} (Equivalence Analysis):} isolates a single claim from \cref{sec:part1}---``IHS
is Dantzig--Wolfe on the dual packing LP''---and asks whether it is literally
true. The answer is yes for an LP master (the $P_4$ example and the exact
equivalence proposition), and no for the integer master used in practice
(the $K_3$ counterexample). The verdict locates the equivalence at the
integrality gap of the hitting-set polytope.

\item \textbf{\Cref{sec:taxonomy} (Broader Taxonomy):} provides a broader taxonomy of the relationships
between IHS and OR methods. It interprets IHS as a discrete form of
Lagrangian relaxation (dual ascent) and a primal cutting-plane method, and
discusses the dual column-generation view in the space of Maximal
Satisfiable Subsets (MSSes).

\item \textbf{\Cref{sec:part3} (Empirical Landscape):} reads the empirical study of
\citet{petrova2025empirical} on IHS for weighted CSPs through the lens of
the preceding sections, mapping their two design axes onto the in-out
separation and Pareto-cut selection techniques predicted by the
decomposition framing.
\end{itemize}

Readers familiar with the core/correction-set duality and with column
generation may skip directly to \cref{sec:part2}; readers interested only in
solver design choices may skip directly to \cref{sec:part3}.
